% Options for packages loaded elsewhere
\PassOptionsToPackage{unicode}{hyperref}
\PassOptionsToPackage{hyphens}{url}
%
\documentclass[
]{article}
\usepackage{amsmath,amssymb}
\usepackage{lmodern}
\usepackage{iftex}
\ifPDFTeX
  \usepackage[T1]{fontenc}
  \usepackage[utf8]{inputenc}
  \usepackage{textcomp} % provide euro and other symbols
\else % if luatex or xetex
  \usepackage{unicode-math}
  \defaultfontfeatures{Scale=MatchLowercase}
  \defaultfontfeatures[\rmfamily]{Ligatures=TeX,Scale=1}
\fi
% Use upquote if available, for straight quotes in verbatim environments
\IfFileExists{upquote.sty}{\usepackage{upquote}}{}
\IfFileExists{microtype.sty}{% use microtype if available
  \usepackage[]{microtype}
  \UseMicrotypeSet[protrusion]{basicmath} % disable protrusion for tt fonts
}{}
\makeatletter
\@ifundefined{KOMAClassName}{% if non-KOMA class
  \IfFileExists{parskip.sty}{%
    \usepackage{parskip}
  }{% else
    \setlength{\parindent}{0pt}
    \setlength{\parskip}{6pt plus 2pt minus 1pt}}
}{% if KOMA class
  \KOMAoptions{parskip=half}}
\makeatother
\usepackage{xcolor}
\usepackage{longtable,booktabs,array}
\usepackage{multirow}
\usepackage{calc} % for calculating minipage widths
% Correct order of tables after \paragraph or \subparagraph
\usepackage{etoolbox}
\makeatletter
\patchcmd\longtable{\par}{\if@noskipsec\mbox{}\fi\par}{}{}
\makeatother
% Allow footnotes in longtable head/foot
\IfFileExists{footnotehyper.sty}{\usepackage{footnotehyper}}{\usepackage{footnote}}
\makesavenoteenv{longtable}
\usepackage{graphicx}
\makeatletter
\def\maxwidth{\ifdim\Gin@nat@width>\linewidth\linewidth\else\Gin@nat@width\fi}
\def\maxheight{\ifdim\Gin@nat@height>\textheight\textheight\else\Gin@nat@height\fi}
\makeatother
% Scale images if necessary, so that they will not overflow the page
% margins by default, and it is still possible to overwrite the defaults
% using explicit options in \includegraphics[width, height, ...]{}
\setkeys{Gin}{width=\maxwidth,height=\maxheight,keepaspectratio}
% Set default figure placement to htbp
\makeatletter
\def\fps@figure{htbp}
\makeatother
\setlength{\emergencystretch}{3em} % prevent overfull lines
\providecommand{\tightlist}{%
  \setlength{\itemsep}{0pt}\setlength{\parskip}{0pt}}
\setcounter{secnumdepth}{-\maxdimen} % remove section numbering
\ifLuaTeX
  \usepackage{selnolig}  % disable illegal ligatures
\fi
\IfFileExists{bookmark.sty}{\usepackage{bookmark}}{\usepackage{hyperref}}
\IfFileExists{xurl.sty}{\usepackage{xurl}}{} % add URL line breaks if available
\urlstyle{same} % disable monospaced font for URLs
\hypersetup{
  hidelinks,
  pdfcreator={LaTeX via pandoc}}

\author{}
\date{}

\begin{document}

\begin{longtable}[]{@{}
  >{\raggedright\arraybackslash}p{(\columnwidth - 2\tabcolsep) * \real{0.5000}}
  >{\raggedright\arraybackslash}p{(\columnwidth - 2\tabcolsep) * \real{0.5000}}@{}}
\toprule()
\begin{minipage}[b]{\linewidth}\raggedright
\begin{quote}
ICS1512
\end{quote}
\end{minipage} & \begin{minipage}[b]{\linewidth}\raggedright
Machine Learning Algorithms Laboratory
\end{minipage} \\
\midrule()
\endhead
\bottomrule()
\end{longtable}

\begin{longtable}[]{@{}
  >{\raggedright\arraybackslash}p{(\columnwidth - 2\tabcolsep) * \real{0.5000}}
  >{\raggedright\arraybackslash}p{(\columnwidth - 2\tabcolsep) * \real{0.5000}}@{}}
\toprule()
\begin{minipage}[b]{\linewidth}\raggedright
\begin{quote}
Experiment No.
\end{quote}
\end{minipage} & \begin{minipage}[b]{\linewidth}\raggedright
\begin{quote}
04
\end{quote}
\end{minipage} \\
\midrule()
\endhead
\begin{minipage}[t]{\linewidth}\raggedright
\begin{quote}
Experiment Title
\end{quote}
\end{minipage} & Binary Classification using Linear and Kernel-Based Models \\
\begin{minipage}[t]{\linewidth}\raggedright
\begin{quote}
GitHub Repository:
\end{quote}
\end{minipage} & \url{https://github.com/The-FOOL-00/ML_Lab} \\
\begin{minipage}[t]{\linewidth}\raggedright
\begin{quote}
Student Name
\end{quote}
\end{minipage} & \begin{minipage}[t]{\linewidth}\raggedright
\begin{quote}
R. Sundareswaran
\end{quote}
\end{minipage} \\
\begin{minipage}[t]{\linewidth}\raggedright
\begin{quote}
Register Number
\end{quote}
\end{minipage} & \begin{minipage}[t]{\linewidth}\raggedright
\begin{quote}
3122247001048
\end{quote}
\end{minipage} \\
\begin{minipage}[t]{\linewidth}\raggedright
\begin{quote}
Faculty
\end{quote}
\end{minipage} & \begin{minipage}[t]{\linewidth}\raggedright
\begin{quote}
Dr. Poreddy Ajay Kumar Reddy
\end{quote}
\end{minipage} \\
\begin{minipage}[t]{\linewidth}\raggedright
\begin{quote}
Submission Date
\end{quote}
\end{minipage} & \begin{minipage}[t]{\linewidth}\raggedright
\begin{quote}
August 11, 2026
\end{quote}
\end{minipage} \\
\bottomrule()
\end{longtable}

\begin{longtable}[]{@{}
  >{\raggedright\arraybackslash}p{(\columnwidth - 2\tabcolsep) * \real{0.5000}}
  >{\raggedright\arraybackslash}p{(\columnwidth - 2\tabcolsep) * \real{0.5000}}@{}}
\toprule()
\begin{minipage}[b]{\linewidth}\raggedright
\begin{quote}
\textbf{1}
\end{quote}
\end{minipage} & \begin{minipage}[b]{\linewidth}\raggedright
\begin{quote}
\textbf{Objective}
\end{quote}
\end{minipage} \\
\midrule()
\endhead
\bottomrule()
\end{longtable}

\begin{quote}
To classify emails as spam or ham using Logistic Regression and Support
Vector Machine (SVM) classifiers and to analyze the effect of
hyperparameter tuning on classification performance.

\textbf{2} \textbf{Problem Statement}

For the \textbf{Spambase} dataset that contains numerical features
extracted from email con-tent and a binary label indicating \textbf{Spam
or non-spam (ham)}. Implement Logistic Regression and reduce overfitting
by \textbf{Regularization}. Tune with Grid Search CV and
RandomizedSearchCV.
\end{quote}

\begin{longtable}[]{@{}
  >{\raggedright\arraybackslash}p{(\columnwidth - 2\tabcolsep) * \real{0.5000}}
  >{\raggedright\arraybackslash}p{(\columnwidth - 2\tabcolsep) * \real{0.5000}}@{}}
\toprule()
\begin{minipage}[b]{\linewidth}\raggedright
\begin{quote}
\textbf{3}
\end{quote}
\end{minipage} & \begin{minipage}[b]{\linewidth}\raggedright
\begin{quote}
\textbf{Dataset Description}
\end{quote}
\end{minipage} \\
\midrule()
\endhead
\bottomrule()
\end{longtable}

\begin{longtable}[]{@{}
  >{\raggedright\arraybackslash}p{(\columnwidth - 2\tabcolsep) * \real{0.5000}}
  >{\raggedright\arraybackslash}p{(\columnwidth - 2\tabcolsep) * \real{0.5000}}@{}}
\toprule()
\begin{minipage}[b]{\linewidth}\raggedright
\begin{quote}
Dataset Name
\end{quote}
\end{minipage} & \begin{minipage}[b]{\linewidth}\raggedright
\begin{quote}
Spambase
\end{quote}
\end{minipage} \\
\midrule()
\endhead
\begin{minipage}[t]{\linewidth}\raggedright
\begin{quote}
Dataset Source
\end{quote}
\end{minipage} & \begin{minipage}[t]{\linewidth}\raggedright
\url{https://www.kaggle.com/datasets/somesh24/spambase}
\end{minipage} \\
\begin{minipage}[t]{\linewidth}\raggedright
\begin{quote}
Number of Samples
\end{quote}
\end{minipage} & \begin{minipage}[t]{\linewidth}\raggedright
\begin{quote}
4601
\end{quote}
\end{minipage} \\
\begin{minipage}[t]{\linewidth}\raggedright
\begin{quote}
Number of Features
\end{quote}
\end{minipage} & \begin{minipage}[t]{\linewidth}\raggedright
\begin{quote}
58
\end{quote}
\end{minipage} \\
\begin{minipage}[t]{\linewidth}\raggedright
\begin{quote}
Number of Classes
\end{quote}
\end{minipage} & \begin{minipage}[t]{\linewidth}\raggedright
\begin{quote}
2
\end{quote}
\end{minipage} \\
\begin{minipage}[t]{\linewidth}\raggedright
\begin{quote}
Missing Values
\end{quote}
\end{minipage} & \begin{minipage}[t]{\linewidth}\raggedright
\begin{quote}
0
\end{quote}
\end{minipage} \\
\begin{minipage}[t]{\linewidth}\raggedright
\begin{quote}
Train-Test Split
\end{quote}
\end{minipage} & \begin{minipage}[t]{\linewidth}\raggedright
\begin{quote}
Train: 3680 and Test: 921
\end{quote}
\end{minipage} \\
\bottomrule()
\end{longtable}

1

\begin{longtable}[]{@{}
  >{\raggedright\arraybackslash}p{(\columnwidth - 4\tabcolsep) * \real{0.3333}}
  >{\raggedright\arraybackslash}p{(\columnwidth - 4\tabcolsep) * \real{0.3333}}
  >{\raggedright\arraybackslash}p{(\columnwidth - 4\tabcolsep) * \real{0.3333}}@{}}
\toprule()
\multicolumn{2}{@{}>{\raggedright\arraybackslash}p{(\columnwidth - 4\tabcolsep) * \real{0.6667} + 2\tabcolsep}}{%
\begin{minipage}[b]{\linewidth}\raggedright
\begin{quote}
ICS1512
\end{quote}
\end{minipage}} & \begin{minipage}[b]{\linewidth}\raggedright
Machine Learning Algorithms Laboratory
\end{minipage} \\
\midrule()
\endhead
\begin{minipage}[t]{\linewidth}\raggedright
\begin{quote}
\textbf{4}
\end{quote}
\end{minipage} &
\multicolumn{2}{>{\raggedright\arraybackslash}p{(\columnwidth - 4\tabcolsep) * \real{0.6667} + 2\tabcolsep}@{}}{%
\begin{minipage}[t]{\linewidth}\raggedright
\begin{quote}
\textbf{Exploratory Data Analysis}
\end{quote}
\end{minipage}} \\
\bottomrule()
\end{longtable}

\begin{quote}
\textbf{Class Distribution}
\end{quote}

\includegraphics[width=4.3875in,height=3.13056in]{media/image1.png}

Figure 1: Class Distribution of Emails

\begin{quote}
\textbf{Inference:}\\
The class distribution plot shows that the dataset contains more ham
emails than spam emails. Approximately 60.6\% of the samples belong to
the ham class, while 39.4\% belong to the spam class. Although the
dataset is not perfectly balanced, the class imbalance is moderate and
suitable for training classification models without extensive
resampling.
\end{quote}

2

\begin{longtable}[]{@{}
  >{\raggedright\arraybackslash}p{(\columnwidth - 2\tabcolsep) * \real{0.5000}}
  >{\raggedright\arraybackslash}p{(\columnwidth - 2\tabcolsep) * \real{0.5000}}@{}}
\toprule()
\begin{minipage}[b]{\linewidth}\raggedright
\begin{quote}
ICS1512
\end{quote}
\end{minipage} & \begin{minipage}[b]{\linewidth}\raggedright
Machine Learning Algorithms Laboratory
\end{minipage} \\
\midrule()
\endhead
\bottomrule()
\end{longtable}

\begin{quote}
\textbf{Correlation Matrix}
\end{quote}

\includegraphics[width=4.3875in,height=3.13333in]{media/image2.png}

Figure 2: Correlation Matrix of Features

\begin{quote}
\textbf{Inference:}\\
The correlation matrix illustrates the relationships among the numerical
features in the Spambase dataset. Most features exhibit weak to moderate
correlations, while a few groups of features show stronger positive
correlations. This indicates that the dataset contains diverse
information with limited redundancy, making it suitable for spam
clas-sification.

\textbf{5} \textbf{Data Preprocessing}

• \textbf{Missing Value Handling:}\\
The Spambase dataset contains no missing values; therefore, no
imputation or removal of records was required.

• \textbf{Encoding:}\\
The target variable is already represented as binary labels (Spam/Ham),
so no additional encoding was necessary.

• \textbf{Feature Scaling:}\\
All numerical features were standardized using StandardScaler. Feature
scaling improves the convergence of Logistic Regression and is essential
for Support Vector Machine since it is distance-based.

• \textbf{Train-Test Split:}\\
The dataset was divided into training and testing sets using an 80:20
ratio, resulting in 3680 training samples and 921 testing samples.

• \textbf{Feature Engineering:}\\
No additional feature engineering techniques were applied since the
dataset already contains extracted numerical features.
\end{quote}

3

\begin{longtable}[]{@{}
  >{\raggedright\arraybackslash}p{(\columnwidth - 4\tabcolsep) * \real{0.3333}}
  >{\raggedright\arraybackslash}p{(\columnwidth - 4\tabcolsep) * \real{0.3333}}
  >{\raggedright\arraybackslash}p{(\columnwidth - 4\tabcolsep) * \real{0.3333}}@{}}
\toprule()
\multicolumn{2}{@{}>{\raggedright\arraybackslash}p{(\columnwidth - 4\tabcolsep) * \real{0.6667} + 2\tabcolsep}}{%
\begin{minipage}[b]{\linewidth}\raggedright
\begin{quote}
ICS1512
\end{quote}
\end{minipage}} & \begin{minipage}[b]{\linewidth}\raggedright
Machine Learning Algorithms Laboratory
\end{minipage} \\
\midrule()
\endhead
\begin{minipage}[t]{\linewidth}\raggedright
\begin{quote}
\textbf{6}
\end{quote}
\end{minipage} &
\multicolumn{2}{>{\raggedright\arraybackslash}p{(\columnwidth - 4\tabcolsep) * \real{0.6667} + 2\tabcolsep}@{}}{%
\begin{minipage}[t]{\linewidth}\raggedright
\begin{quote}
\textbf{Machine Learning Algorithm}
\end{quote}
\end{minipage}} \\
\textbf{6.1} &
\multicolumn{2}{>{\raggedright\arraybackslash}p{(\columnwidth - 4\tabcolsep) * \real{0.6667} + 2\tabcolsep}@{}}{%
\begin{minipage}[t]{\linewidth}\raggedright
\begin{quote}
\textbf{Working Principle}
\end{quote}
\end{minipage}} \\
\bottomrule()
\end{longtable}

\begin{quote}
• \textbf{Logistic Regression:}\\
Logistic Regression is a probabilistic classification algorithm used for
binary classi-fication problems. It models the probability that a sample
belongs to a particular class using the \textbf{sigmoid function}.

• \textbf{Support Vector Machine (SVM)} Support Vector Machine is a
margin-based classifier that finds an optimal hyperplane seperating two
classes by maximizing the margin between them. \textbf{Key Idea:} Only a
subset of training points (Support vectors) define the decision
boundary.

\textbf{Mathematical Formulation}

\[P(y = 1|x) = \frac{1}{1 + e^{-(w^T x + b)}}\]

A threshold (usually 0.5) is applied to convert probability into class
labels.

\textbf{Loss Function}: Logistic Regression minimizes the
\textbf{log-loss (cross-entropy loss)}, which penalizes incorrect
predictions.

\textbf{6.2} \textbf{Advantages}

• Fast and Simple.

• Works well with linearly seperable data.

• \textbf{Less prone to overfitting:} with L1/L2 Regularization.

• Produces probability estimates for binary classification.

\textbf{6.3} \textbf{Limitations}

• Assumes Linear Decision boundary.

• Sensitive to outliers.

• Sensitive to multicollinearity.

• Poor performance if classes are not linearly seperable.
\end{quote}

\begin{longtable}[]{@{}
  >{\raggedright\arraybackslash}p{(\columnwidth - 4\tabcolsep) * \real{0.3333}}
  >{\raggedright\arraybackslash}p{(\columnwidth - 4\tabcolsep) * \real{0.3333}}
  >{\raggedright\arraybackslash}p{(\columnwidth - 4\tabcolsep) * \real{0.3333}}@{}}
\toprule()
\begin{minipage}[b]{\linewidth}\raggedright
\begin{quote}
\textbf{7}
\end{quote}
\end{minipage} &
\multicolumn{2}{>{\raggedright\arraybackslash}p{(\columnwidth - 4\tabcolsep) * \real{0.6667} + 2\tabcolsep}@{}}{%
\begin{minipage}[b]{\linewidth}\raggedright
\begin{quote}
\textbf{Experimental Setup}
\end{quote}
\end{minipage}} \\
\midrule()
\endhead
\multicolumn{2}{@{}>{\raggedright\arraybackslash}p{(\columnwidth - 4\tabcolsep) * \real{0.6667} + 2\tabcolsep}}{%
\begin{minipage}[t]{\linewidth}\raggedright
\begin{quote}
Python Version\\
Scikit-Learn Version IDE\\
Operating System\\
Hardware
\end{quote}\strut
\end{minipage}} & \begin{minipage}[t]{\linewidth}\raggedright
\begin{quote}
3.13.13\\
1.7.5\\
Jupyter Notebook\\
Fedora 42\\
Intel Core i5 (11th gen), 8GB RAM
\end{quote}\strut
\end{minipage} \\
\bottomrule()
\end{longtable}

4

\begin{longtable}[]{@{}
  >{\raggedright\arraybackslash}p{(\columnwidth - 4\tabcolsep) * \real{0.3333}}
  >{\raggedright\arraybackslash}p{(\columnwidth - 4\tabcolsep) * \real{0.3333}}
  >{\raggedright\arraybackslash}p{(\columnwidth - 4\tabcolsep) * \real{0.3333}}@{}}
\toprule()
\multicolumn{2}{@{}>{\raggedright\arraybackslash}p{(\columnwidth - 4\tabcolsep) * \real{0.6667} + 2\tabcolsep}}{%
\begin{minipage}[b]{\linewidth}\raggedright
\begin{quote}
ICS1512
\end{quote}
\end{minipage}} & \begin{minipage}[b]{\linewidth}\raggedright
Machine Learning Algorithms Laboratory
\end{minipage} \\
\midrule()
\endhead
\begin{minipage}[t]{\linewidth}\raggedright
\begin{quote}
\textbf{8}
\end{quote}
\end{minipage} &
\multicolumn{2}{>{\raggedright\arraybackslash}p{(\columnwidth - 4\tabcolsep) * \real{0.6667} + 2\tabcolsep}@{}}{%
\begin{minipage}[t]{\linewidth}\raggedright
\begin{quote}
\textbf{Hyperparameter Selection}
\end{quote}
\end{minipage}} \\
\textbf{8.1} &
\multicolumn{2}{>{\raggedright\arraybackslash}p{(\columnwidth - 4\tabcolsep) * \real{0.6667} + 2\tabcolsep}@{}}{%
\begin{minipage}[t]{\linewidth}\raggedright
\begin{quote}
\textbf{GridSearchCV for LogisticRegression}
\end{quote}
\end{minipage}} \\
\multicolumn{2}{@{}>{\raggedright\arraybackslash}p{(\columnwidth - 4\tabcolsep) * \real{0.6667} + 2\tabcolsep}}{%
Parameter} & \begin{minipage}[t]{\linewidth}\raggedright
\begin{quote}
Value
\end{quote}
\end{minipage} \\
\multicolumn{2}{@{}>{\raggedright\arraybackslash}p{(\columnwidth - 4\tabcolsep) * \real{0.6667} + 2\tabcolsep}}{%
\begin{minipage}[t]{\linewidth}\raggedright
\begin{quote}
C\\
solver
\end{quote}\strut
\end{minipage}} & \begin{minipage}[t]{\linewidth}\raggedright
\begin{quote}
10\\
liblinear
\end{quote}\strut
\end{minipage} \\
\bottomrule()
\end{longtable}

\begin{quote}
\textbf{GridSearchCV for SVM}
\end{quote}

\begin{longtable}[]{@{}
  >{\raggedright\arraybackslash}p{(\columnwidth - 14\tabcolsep) * \real{0.1250}}
  >{\raggedright\arraybackslash}p{(\columnwidth - 14\tabcolsep) * \real{0.1250}}
  >{\raggedright\arraybackslash}p{(\columnwidth - 14\tabcolsep) * \real{0.1250}}
  >{\raggedright\arraybackslash}p{(\columnwidth - 14\tabcolsep) * \real{0.1250}}
  >{\raggedright\arraybackslash}p{(\columnwidth - 14\tabcolsep) * \real{0.1250}}
  >{\raggedright\arraybackslash}p{(\columnwidth - 14\tabcolsep) * \real{0.1250}}
  >{\raggedright\arraybackslash}p{(\columnwidth - 14\tabcolsep) * \real{0.1250}}
  >{\raggedright\arraybackslash}p{(\columnwidth - 14\tabcolsep) * \real{0.1250}}@{}}
\toprule()
\multicolumn{2}{@{}>{\raggedright\arraybackslash}p{(\columnwidth - 14\tabcolsep) * \real{0.2500} + 2\tabcolsep}}{%
\begin{minipage}[b]{\linewidth}\raggedright
Parameter
\end{minipage}} & \begin{minipage}[b]{\linewidth}\raggedright
\begin{quote}
Value
\end{quote}
\end{minipage} & \begin{minipage}[b]{\linewidth}\raggedright
\end{minipage} & \begin{minipage}[b]{\linewidth}\raggedright
\end{minipage} & \begin{minipage}[b]{\linewidth}\raggedright
\end{minipage} & \begin{minipage}[b]{\linewidth}\raggedright
\end{minipage} & \begin{minipage}[b]{\linewidth}\raggedright
\end{minipage} \\
\midrule()
\endhead
\multicolumn{2}{@{}>{\raggedright\arraybackslash}p{(\columnwidth - 14\tabcolsep) * \real{0.2500} + 2\tabcolsep}}{%
\begin{minipage}[t]{\linewidth}\raggedright
\begin{quote}
C\\
gamma\\
kernel\\
degree
\end{quote}\strut
\end{minipage}} & \begin{minipage}[t]{\linewidth}\raggedright
\begin{quote}
1\\
scale\\
RBF\\
2
\end{quote}\strut
\end{minipage} & & & & & \\
\begin{minipage}[t]{\linewidth}\raggedright
\begin{quote}
\textbf{9}
\end{quote}
\end{minipage} &
\multicolumn{2}{>{\raggedright\arraybackslash}p{(\columnwidth - 14\tabcolsep) * \real{0.2500} + 2\tabcolsep}}{%
\begin{minipage}[t]{\linewidth}\raggedright
\begin{quote}
\textbf{Results}
\end{quote}
\end{minipage}} & & & & & \\
\multicolumn{3}{@{}>{\raggedright\arraybackslash}p{(\columnwidth - 14\tabcolsep) * \real{0.3750} + 4\tabcolsep}}{%
\begin{minipage}[t]{\linewidth}\raggedright
\begin{quote}
Model
\end{quote}
\end{minipage}} & Accuracy & Precision & Recall & F1-Score & ROC-AUC \\
\multicolumn{3}{@{}>{\raggedright\arraybackslash}p{(\columnwidth - 14\tabcolsep) * \real{0.3750} + 4\tabcolsep}}{%
\begin{minipage}[t]{\linewidth}\raggedright
\begin{quote}
LogisticRegression\\
Support Vector Machine
\end{quote}\strut
\end{minipage}} & \begin{minipage}[t]{\linewidth}\raggedright
\begin{quote}
0.9196\\
0.925
\end{quote}
\end{minipage} & \begin{minipage}[t]{\linewidth}\raggedright
\begin{quote}
0.9316\\
0.934
\end{quote}
\end{minipage} & \begin{minipage}[t]{\linewidth}\raggedright
0.8743\\
0.8846
\end{minipage} &
\begin{minipage}[t]{\linewidth}\raggedright
\begin{quote}
09021\\
0.9090
\end{quote}\strut
\end{minipage} & \begin{minipage}[t]{\linewidth}\raggedright
\begin{quote}
0.97\\
1.00
\end{quote}\strut
\end{minipage} \\
\bottomrule()
\end{longtable}

\begin{quote}
Table 2: Comparing performance of LogisticRegression and SVM .
\end{quote}

\begin{longtable}[]{@{}
  >{\raggedright\arraybackslash}p{(\columnwidth - 4\tabcolsep) * \real{0.3333}}
  >{\raggedright\arraybackslash}p{(\columnwidth - 4\tabcolsep) * \real{0.3333}}
  >{\raggedright\arraybackslash}p{(\columnwidth - 4\tabcolsep) * \real{0.3333}}@{}}
\toprule()
\begin{minipage}[b]{\linewidth}\raggedright
\begin{quote}
Kernel
\end{quote}
\end{minipage} & \begin{minipage}[b]{\linewidth}\raggedright
accuracy
\end{minipage} & \begin{minipage}[b]{\linewidth}\raggedright
F1
\end{minipage} \\
\midrule()
\endhead
\begin{minipage}[t]{\linewidth}\raggedright
\begin{quote}
Linear\\
Poly\\
RBF\\
Sigmoid
\end{quote}\strut
\end{minipage} & \begin{minipage}[t]{\linewidth}\raggedright
\begin{quote}
0.925\\
0.764\\
0.934\\
0.889
\end{quote}\strut
\end{minipage} & \begin{minipage}[t]{\linewidth}\raggedright
0.9090\\
0.629\\
0.9206\\
0.8664\strut
\end{minipage} \\
\bottomrule()
\end{longtable}

Table 3: Comparing the performance of various kernels of SVC

\begin{longtable}[]{@{}
  >{\raggedright\arraybackslash}p{(\columnwidth - 4\tabcolsep) * \real{0.3333}}
  >{\raggedright\arraybackslash}p{(\columnwidth - 4\tabcolsep) * \real{0.3333}}
  >{\raggedright\arraybackslash}p{(\columnwidth - 4\tabcolsep) * \real{0.3333}}@{}}
\toprule()
\begin{minipage}[b]{\linewidth}\raggedright
\begin{quote}
Model
\end{quote}
\end{minipage} & \begin{minipage}[b]{\linewidth}\raggedright
Train accuracy
\end{minipage} & \begin{minipage}[b]{\linewidth}\raggedright
Test accuracy
\end{minipage} \\
\midrule()
\endhead
\begin{minipage}[t]{\linewidth}\raggedright
\begin{quote}
Logistic regression\\
SVM
\end{quote}
\end{minipage} & \begin{minipage}[t]{\linewidth}\raggedright
\begin{quote}
0.9277\\
0.945
\end{quote}\strut
\end{minipage} & \begin{minipage}[t]{\linewidth}\raggedright
\begin{quote}
0.9196\\
0.9348
\end{quote}\strut
\end{minipage} \\
\bottomrule()
\end{longtable}

Table 4: Comparing Test and train accuracies of both classification
models.

5

\begin{longtable}[]{@{}
  >{\raggedright\arraybackslash}p{(\columnwidth - 4\tabcolsep) * \real{0.3333}}
  >{\raggedright\arraybackslash}p{(\columnwidth - 4\tabcolsep) * \real{0.3333}}
  >{\raggedright\arraybackslash}p{(\columnwidth - 4\tabcolsep) * \real{0.3333}}@{}}
\toprule()
\multicolumn{2}{@{}>{\raggedright\arraybackslash}p{(\columnwidth - 4\tabcolsep) * \real{0.6667} + 2\tabcolsep}}{%
\begin{minipage}[b]{\linewidth}\raggedright
\begin{quote}
ICS1512
\end{quote}
\end{minipage}} & \begin{minipage}[b]{\linewidth}\raggedright
\begin{quote}
Machine Learning Algorithms Laboratory
\end{quote}
\end{minipage} \\
\midrule()
\endhead
\begin{minipage}[t]{\linewidth}\raggedright
\begin{quote}
\textbf{10}
\end{quote}
\end{minipage} & \begin{minipage}[t]{\linewidth}\raggedright
\begin{quote}
\textbf{Result Visualization}
\end{quote}
\end{minipage} & \multirow{2}{*}{} \\
\textbf{10.1} & \begin{minipage}[t]{\linewidth}\raggedright
\begin{quote}
\textbf{Logistic Regression}
\end{quote}
\end{minipage} \\
\bottomrule()
\end{longtable}

\includegraphics[width=5.01389in,height=4.28194in]{media/image3.png}

Figure 3: Confusion Matrix

6

\begin{longtable}[]{@{}
  >{\raggedright\arraybackslash}p{(\columnwidth - 2\tabcolsep) * \real{0.5000}}
  >{\raggedright\arraybackslash}p{(\columnwidth - 2\tabcolsep) * \real{0.5000}}@{}}
\toprule()
\begin{minipage}[b]{\linewidth}\raggedright
\begin{quote}
ICS1512
\end{quote}
\end{minipage} & \begin{minipage}[b]{\linewidth}\raggedright
Machine Learning Algorithms Laboratory
\end{minipage} \\
\midrule()
\endhead
\bottomrule()
\end{longtable}

\includegraphics[width=5.01389in,height=4.27222in]{media/image4.png}

Figure 4: Confusion Matrix Tuned with GridSearchCV.

7

\begin{longtable}[]{@{}
  >{\raggedright\arraybackslash}p{(\columnwidth - 2\tabcolsep) * \real{0.5000}}
  >{\raggedright\arraybackslash}p{(\columnwidth - 2\tabcolsep) * \real{0.5000}}@{}}
\toprule()
\begin{minipage}[b]{\linewidth}\raggedright
\begin{quote}
ICS1512
\end{quote}
\end{minipage} & \begin{minipage}[b]{\linewidth}\raggedright
Machine Learning Algorithms Laboratory
\end{minipage} \\
\midrule()
\endhead
\bottomrule()
\end{longtable}

\includegraphics[width=5.01389in,height=4.85833in]{media/image5.png}

Figure 5: ROC curve

8

\begin{longtable}[]{@{}
  >{\raggedright\arraybackslash}p{(\columnwidth - 2\tabcolsep) * \real{0.5000}}
  >{\raggedright\arraybackslash}p{(\columnwidth - 2\tabcolsep) * \real{0.5000}}@{}}
\toprule()
\begin{minipage}[b]{\linewidth}\raggedright
\begin{quote}
ICS1512
\end{quote}
\end{minipage} & \begin{minipage}[b]{\linewidth}\raggedright
Machine Learning Algorithms Laboratory
\end{minipage} \\
\midrule()
\endhead
\bottomrule()
\end{longtable}

\includegraphics[width=5.01389in,height=4.85833in]{media/image6.png}

Figure 6: ROC curve Tuned with GridSearchCV.

9

\begin{longtable}[]{@{}
  >{\raggedright\arraybackslash}p{(\columnwidth - 4\tabcolsep) * \real{0.3333}}
  >{\raggedright\arraybackslash}p{(\columnwidth - 4\tabcolsep) * \real{0.3333}}
  >{\raggedright\arraybackslash}p{(\columnwidth - 4\tabcolsep) * \real{0.3333}}@{}}
\toprule()
\begin{minipage}[b]{\linewidth}\raggedright
ICS1512
\end{minipage} & \begin{minipage}[b]{\linewidth}\raggedright
\end{minipage} & \begin{minipage}[b]{\linewidth}\raggedright
Machine Learning Algorithms Laboratory
\end{minipage} \\
\midrule()
\endhead
\begin{minipage}[t]{\linewidth}\raggedright
\begin{quote}
\textbf{10.2}
\end{quote}
\end{minipage} & \begin{minipage}[t]{\linewidth}\raggedright
\begin{quote}
\textbf{SVM}
\end{quote}
\end{minipage} & \\
\bottomrule()
\end{longtable}

\includegraphics[width=5.01389in,height=4.27222in]{media/image7.png}

Figure 7: Confusion Matrix

10

\begin{longtable}[]{@{}
  >{\raggedright\arraybackslash}p{(\columnwidth - 2\tabcolsep) * \real{0.5000}}
  >{\raggedright\arraybackslash}p{(\columnwidth - 2\tabcolsep) * \real{0.5000}}@{}}
\toprule()
\begin{minipage}[b]{\linewidth}\raggedright
\begin{quote}
ICS1512
\end{quote}
\end{minipage} & \begin{minipage}[b]{\linewidth}\raggedright
Machine Learning Algorithms Laboratory
\end{minipage} \\
\midrule()
\endhead
\bottomrule()
\end{longtable}

\includegraphics[width=5.01389in,height=4.27222in]{media/image8.png}

Figure 8: Confusion Matrix Tuned with GridSearchCV.

11

\begin{longtable}[]{@{}
  >{\raggedright\arraybackslash}p{(\columnwidth - 2\tabcolsep) * \real{0.5000}}
  >{\raggedright\arraybackslash}p{(\columnwidth - 2\tabcolsep) * \real{0.5000}}@{}}
\toprule()
\begin{minipage}[b]{\linewidth}\raggedright
\begin{quote}
ICS1512
\end{quote}
\end{minipage} & \begin{minipage}[b]{\linewidth}\raggedright
Machine Learning Algorithms Laboratory
\end{minipage} \\
\midrule()
\endhead
\bottomrule()
\end{longtable}

\includegraphics[width=5.01389in,height=4.85833in]{media/image9.png}

Figure 9: ROC curve

12

\begin{longtable}[]{@{}
  >{\raggedright\arraybackslash}p{(\columnwidth - 2\tabcolsep) * \real{0.5000}}
  >{\raggedright\arraybackslash}p{(\columnwidth - 2\tabcolsep) * \real{0.5000}}@{}}
\toprule()
\begin{minipage}[b]{\linewidth}\raggedright
\begin{quote}
ICS1512
\end{quote}
\end{minipage} & \begin{minipage}[b]{\linewidth}\raggedright
Machine Learning Algorithms Laboratory
\end{minipage} \\
\midrule()
\endhead
\bottomrule()
\end{longtable}

\includegraphics[width=5.01389in,height=4.85833in]{media/image10.png}

Figure 10: ROC curve tuned with gridSearchCV.

\begin{quote}
\textbf{11 Discussion}

Since I have trained, GridSearchCV for LogRegression with max\_iter=10000
it takes nearly 1min 10secs to finish. Also GridSearchCV for SVM
with max\_iter=5000 that takes 2mins roughly to finish. So these time
complexities are the limitations of my implementation.

\textbf{12 Conclusion}

From this experiment, I have learnt to implement Logistic Regression and
Support Vector Classification for a given dataset having binary class
target. Also tuned them using GridSearchCV, plotted Confusion matrix and
ROC curve for all 4 models. Compared Performance metrics for all 4
kernels of SVM.

\textbf{13 References}

1. Scikit-learn Developers, ``Scikit-learn: Machine Learning in Python.''
Available: \url{https://scikit-learn.org/stable/}
\end{quote}

13

\begin{longtable}[]{@{}
  >{\raggedright\arraybackslash}p{(\columnwidth - 2\tabcolsep) * \real{0.5000}}
  >{\raggedright\arraybackslash}p{(\columnwidth - 2\tabcolsep) * \real{0.5000}}@{}}
\toprule()
\begin{minipage}[b]{\linewidth}\raggedright
\begin{quote}
ICS1512
\end{quote}
\end{minipage} & \begin{minipage}[b]{\linewidth}\raggedright
Machine Learning Algorithms Laboratory
\end{minipage} \\
\midrule()
\endhead
\bottomrule()
\end{longtable}

\begin{quote}
2. Kaggle, ``Spambase Dataset.'' Available:
\url{https://www.kaggle.com/datasets/somesh24/spambase}

3. T. Hastie, R. Tibshirani and J. Friedman, \emph{The Elements of
Statistical Learning}, 2nd ed., Springer, 2009.

4. C. M. Bishop, \emph{Pattern Recognition and Machine Learning},
Springer, 2006.

5. A. Géron, \emph{Hands-On Machine Learning with Scikit-Learn, Keras
and TensorFlow}, 3rd ed., O'Reilly, 2022.
\end{quote}

14

\end{document}
